\section{Related literature}
\label{sec:literature}

\paragraph{Automation in growth models.} The mechanism simulated here is the oldest one in the automation literature. \citet{zeira1998} shows that when technology moves tasks from labour to capital, output rises while labour's share of it falls, and the aggregate capital exponent rises with the automated share of tasks. \citet{acemoglu2018} add a countervailing force, separating the displacement effect of automation from the reinstatement effect of new labour-intensive tasks. \citet{acemoglu2022tasks} find that task displacement accounts for a substantial part of the post-1980 rise in US wage inequality. In that framework the exponent is an equilibrium object. Here it is a parameter, which is the cost of running the mechanism inside a calibrated fiscal overlapping-generations model.

\paragraph{Automation, wealth and factor shares.} \citet{moll2022} show that automation redistributes towards wealth: capital income is more concentrated than labour income, so a shock that raises the return on wealth and lowers the labour share raises inequality through portfolios as well as wages. \citet{mollimas2026} raise the capital exponent from $1/3$ towards~1 and argue for an aggregate elasticity of substitution well below one. % TODO before submission: check the elasticity attributions to mollimas2026 and korinek2026 against the PDFs (referee item 6 on #24)
\citet{acemoglu2025} argues that task-level exposure limits the attainable aggregate gain. \citet{jonestonetti2026} and \citet{acemoglu2026} are further contributions to this debate. % TODO verify characterisation of jonestonetti2026 and acemoglu2026 against the papers

\paragraph{Automation in overlapping-generations models.} \citet{sachs2015robots} build an overlapping-generations model of automation in which the generation that owns the machines gains and later generations can lose, because a lower labour share reduces the wage income from which the young save. This paper extends that question to two scenarios built from the published 2026 ones, in a model calibrated to UK data with an estimated tax system, and sets out incidence by birth cohort and lifetime-earnings type together rather than young against old. % TODO verify characterisation of sachs2015robots against the paper

\paragraph{Quantified AI scenarios.} \citet{korinek2026} supply the sizes used here. Their framework has an automation share, an affected task mass, an adoption parameter and a reinstatement parameter that build from a mid-2026 anchor to 2030, and a matching block that produces equilibrium unemployment. The ``substantial'' column of their Table~3 reports, for the US in 2030, a \genKorinekTableLabourShareFallPp{} percentage-point fall in the labour share, measured TFP \genKorinekTableTfpPct{} per cent above the no-AI path, GDP \genKorinekTableGdpPct{} per cent above it, a capital stock \genKorinekTableCapitalPct{} per cent higher and unemployment of \genKorinekTableUnemploymentPct{} per cent. Four of their six AI channels have no counterpart in OG-UK (Appendix~\ref{app:channels}).

\paragraph{The UK official benchmarks.} In the OBR's March~2026 technological-displacement scenario, capital deepens, the level of GDP is unchanged against the central forecast, higher average labour productivity does not pass through to real earnings, the labour share falls and the profit share rises. Equilibrium unemployment rises to 5.5~per cent, and receipts fall because labour income faces a higher effective tax rate than capital income \citep[Box~2.2 and para~6.18]{obr2026efo}. The OBR's central forecast already includes an AI productivity assumption, built up in Annex~B of \citet{obr2025bp9}. The OBR's own overlapping-generations model \citep{obr2025wp22} is close in architecture to OG-UK. It runs either closed or as a small open economy, and has not been used for an AI scenario.

\paragraph{Capital mobility.} Whether extra investment is financed at home or abroad is an old empirical question. \citet{feldsteinhorioka1980} found that national saving and investment moved almost one for one across OECD countries in 1960--74. \citet{blanchardgiavazzi2002} find much lower coefficients for the European Union and the euro area in 1991--2001. Section~\ref{sec:closure} uses these estimates to judge the model's capital closure.

\paragraph{The modelling platform.} OG-Core \citep{debacker2024ogcore} is an open-source descendant of the Auerbach--Kotlikoff framework \citep{auerbach1987}. OG-UK is its UK calibration, with demographics from UN projections and tax functions from PolicyEngine UK.
